\documentclass[12pt]{article}

\usepackage[T1]{fontenc}
\usepackage[utf8]{inputenc}
\usepackage[margin=1in]{geometry}
\usepackage{setspace}
\usepackage{parskip}
\usepackage{titlesec}
\usepackage[hidelinks]{hyperref}

\titleformat{\section}{\normalfont\bfseries\large}{}{0pt}{}
\titlespacing*{\section}{0pt}{1.4em}{0.6em}

\setstretch{1.15}

\begin{document}

\begin{center}
	{\Large\bfseries Primary Source Worksheet \#3}\\[0.6em]
	RELI 252 \quad\textbullet\quad Dr.\ Clark\\[0.3em]
	Greyson Knippel
\end{center}

\vspace{1em}
\hrule
\vspace{1.5em}

\noindent\textbf{Primary Source Title:} \textit{The Confessions of Nat Turner, the Leader of the late Insurrection in Southampton, Va.}

\vspace{0.5em}

\noindent\textbf{Primary Source Author:} Nat Turner, recorded and published by Thomas Gray 

\section*{What is the author's intent with this piece?}

This source has two authors with two different goals, and that is the most important thing about it. Turner is speaking, but Gray is the one writing it down, choosing what to keep, and selling it.

Gray wants the white readers to see Turner as a weird religious fanatic. He asks Turner directly if there was any bigger plan, and writes the answer no. He inserts his own comments in brackets, like when he says the marks on Turner's body were just common bumps. He ends by calling Turner a ``complete fanatic'' with a ``fiend-like face.'' If the rebellion was just one crazy man hearing voices, then slavery itself is not the problem.

Gray also wants to sell pamphlets. He writes that his blood ``curdled in my veins,'' which is not exactly the right language to use in a court record.

Turner's intent is different and it still comes through. He wants to explain that he was called by God, not acting on anger. He starts before his own birth to show the call was there from the beginning. He points out he never stole, never drank, never swore. He is building a case that he was a prophet.

The most interesting moment to me is when Gray asks if he was mistaken. Turner answers, ``Was not Christ crucified?'' He is saying that being executed does not prove he was wrong. It actuallyproves the opposite.

\section*{Compare and contrast with a previously assigned primary source}

Turner's confession and Frederick Douglass's story both argue that Christianity and slavery cannot coexist, but they get to completely different conclusions about what to do about it.

Both men were enslaved, both learned to read, and both treated the Bible as the real authority. Douglass taught over forty people to read on Sundays. Turner taught himself so young he says he never remembers learning the alphabet. and both of them were punished for that literacy. Douglass's school was broken up with sticks and stones by church class leaders. Turner's reading fed the visions that got him hanged.

The big difference though, is what the Bible tells them to do. Douglass separates ``the Christianity of Christ'' from slaveholding religion and uses words to attack the second one. Turner reads the same Bible and finds the end of the world in it. He sees blood on the corn, figures in the sky, and a voice telling him to ``slay my enemies with their own weapons.'' His weapon is an axe, not preaching.

What explains the difference is partly timing and partly place. Turner is in rural Virginia in 1831 with no path out and no audience. Douglass escapes in 1838 and writes from Massachusetts with abolitionists behind him. Douglass had somewhere to put his anger, and Turner did not.

There is also a difference in how we get their words. Douglass wrote his own book. Turner's was written by a white lawyer who wanted him to look insane. So one is a voice, and the other is a voice filtered through someone trying to control it.

Together they show that reading the Bible for yourself was the dangerous part. Whether it produced a newspaper or a rebellion, white authorities treated it as a threat, which is why South Carolina banned night worship right after Turner died.

\end{document}